\subsection*{Abelian Groups}
\label{subsec:abelian-groups}

\begin{definitionbox}{Definition (Abelian Group)}
\begin{definition}[Abelian Group]
\label{def:algebraic-abelian-group}
A group $(G,\star)$ is \textbf{abelian} if its operation is commutative:
\[
\forall a,b \in G,\quad a\star b=b\star a.
\]
\end{definition}
\end{definitionbox}

\begin{remark*}[Standard quantified statement]
\[
\operatorname{AbelianGroup}(G,\star)
\Longleftrightarrow
\operatorname{Group}(G,\star)
\land
\operatorname{Commutative}(\star,G).
\]
\end{remark*}

\begin{remark*}[Definition predicate reading]
$(G,\star)$ is abelian exactly when it is a group whose operation commutes.
\end{remark*}

\begin{remark*}[Negated quantified statement]
\[
\neg\operatorname{AbelianGroup}(G,\star)
\Longleftrightarrow
\neg\operatorname{Group}(G,\star)
\lor
\exists a,b\in G,\; a\star b\neq b\star a.
\]
\end{remark*}

\begin{remark*}[Interpretation]
Abelian groups are usually written additively: the operation is $+$, the
identity is $0$, and the inverse of $a$ is $-a$.
\end{remark*}

\begin{dependencies}
\begin{itemize}
  \item \hyperref[def:algebraic-group]{Group}
  \item \hyperref[def:algebraic-commutativity]{Commutativity}
\end{itemize}
\end{dependencies}
